\documentclass[12pt,a4paper]{article}

% ---------------------------------------------------------------
% Paquetes
% ---------------------------------------------------------------
\usepackage[utf8]{inputenc}
\usepackage[T1]{fontenc}
\usepackage[spanish,es-lcroman]{babel}
\usepackage{lmodern}
\usepackage{microtype}

\usepackage{amsmath,amssymb,amsthm}
\usepackage{mathtools}
\usepackage{bm}

\usepackage{geometry}
\geometry{margin=2.5cm, top=3cm, bottom=3cm}

\usepackage{booktabs}
\usepackage{array}
\usepackage{multirow}
\usepackage{tabularx}

\usepackage{hyperref}
\hypersetup{
  colorlinks=true,
  linkcolor=black,
  citecolor=black,
  urlcolor=blue,
  pdftitle={Análisis de Seguridad Extendido del Cifrado Alaniz},
  pdfauthor={Alfredo Rodríguez Langa},
}

\usepackage{xcolor}
\usepackage{fancyhdr}
\usepackage{titlesec}
\usepackage{abstract}

\usepackage{listings}
\lstset{basicstyle=\ttfamily\small, frame=single, breaklines=true}

% ---------------------------------------------------------------
% Entornos de teoremas
% ---------------------------------------------------------------
\newtheorem{theorem}{Teorema}[section]
\newtheorem{lemma}[theorem]{Lema}
\newtheorem{corollary}[theorem]{Corolario}
\newtheorem{proposition}[theorem]{Proposición}
\theoremstyle{definition}
\newtheorem{definition}[theorem]{Definición}
\newtheorem{remark}[theorem]{Observación}

% ---------------------------------------------------------------
% Macros matemáticas
% ---------------------------------------------------------------
\newcommand{\Fp}{\mathbb{F}_p}
\newcommand{\Fq}{\mathbb{F}_q}
\newcommand{\GL}{\mathrm{GL}}
\newcommand{\Hzero}{H^{0}}
\newcommand{\Img}{\mathrm{Im}}
\newcommand{\Tr}{\mathrm{Tr}}
\newcommand{\bsigma}{\bm{\sigma}}
\newcommand{\calF}{\mathcal{F}}

% ---------------------------------------------------------------
% Encabezado y pie de página
% ---------------------------------------------------------------
\pagestyle{fancy}
\fancyhf{}
\rhead{\small Análisis de Seguridad Extendido del Cifrado Alaniz}
\lhead{\small A.~Rodr\'iguez Langa, INECO}
\cfoot{\thepage}
\renewcommand{\headrulewidth}{0.4pt}

% ---------------------------------------------------------------
% Título
% ---------------------------------------------------------------
\title{%
  \textbf{Análisis de Seguridad Extendido del Cifrado Alaniz}\\[6pt]
  \large Ataques de Texto en Claro Elegido, Recuperación de Clave y\\
  Propiedades Estructurales de la Función de Cifrado Local\\[4pt]
  \normalsize\textit{Análisis complementario a: L.~Alaniz Pintos,
  ``The Alaniz Cipher: Encryption from Nonlinear Sheaf Morphisms
  over Graphs'', 2026}
}

\author{%
  Alfredo Rodr\'iguez Langa\\[4pt]
  \normalsize Gerencia de Área de Smart Products, INECO\\
  \normalsize Madrid, España\\[2pt]
  \small \href{mailto:alfredo.rodriguez@ineco.com}{alfredo.rodriguez@ineco.com}
}

\date{Marzo de 2026}

% ---------------------------------------------------------------
% Documento
% ---------------------------------------------------------------
\begin{document}

\maketitle
\thispagestyle{fancy}

% ---------------------------------------------------------------
\begin{abstract}
Este documento presenta siete experimentos de seguridad reproducibles
(Experimentos~07--13) que complementan y someten a prueba de estrés
las afirmaciones de seguridad del Cifrado Alaniz (SHEC). Los experimentos
analizan la función de cifrado local $g_B(y) = y + B\,\sigma(y)$
en busca de propiedades de colisión, intentan la recuperación de clave
bajo modelos de texto en claro conocido y elegido, y estudian la
estructura algebraica explotable por un adversario.

El hallazgo central se encuentra en el \textbf{Experimento~10}: ambas
funciones no lineales soportadas ($\sigma = x^3$ y $\sigma = x^{p-2}$)
satisfacen una ley de escala componente a componente
$\sigma(\lambda y)_i = \lambda^e\,\sigma(y)_i$, que permite a un
atacante de texto en claro elegido (CPA, por sus siglas en inglés)
recuperar las matrices de clave por nodo $(A_v, B_v)$ usando únicamente
$O(d)$ consultas de cifrado, con independencia del tamaño del
cuerpo~$p$ y del tamaño de $\GL(d,\Fp)$.
Este resultado contradice directamente la afirmación de resistencia CPA
del artículo original y no está cubierto por la reducción de dureza
NL-SMIP, que concierne a un modelo adversarial distinto.
Los experimentos restantes caracterizan la estructura de colisiones de
$g_B$ (Experimentos~07 y~11), miden la complejidad concreta de la
recuperación de clave exhaustiva y diferencial bajo el modelo de texto
en claro conocido (Experimentos~08, 09, 12 y~13), y aportan la base
algebraica necesaria para evaluar cualquier corrección propuesta.
\end{abstract}

\tableofcontents
\newpage

% ================================================================
\section{Introducción}
% ================================================================

El Cifrado Alaniz cifra un mensaje codificado como sección global
$s \in \Hzero(G, \calF)$ de un haz celular sobre un grafo~$G$.
En cada nodo $v$, el texto cifrado local es
\begin{equation}\label{eq:cifrado}
  c_v = A_v \cdot s_v + B_v \cdot \sigma(A_v \cdot s_v),
  \qquad (A_v, B_v) \in \GL(d, \Fp)^2,
\end{equation}
donde $\sigma : \Fp^d \to \Fp^d$ es una función no lineal pública
aplicada componente a componente. El artículo original argumenta que
esta construcción resiste los ataques de texto en claro elegido (CPA)
porque $\sigma$ impide la linealización: cualquier intento de modelar
$c_v$ como función lineal de $s_v$ produce un sistema inconsistente
una vez acumulados suficientes pares de texto cifrado.

Los experimentos de este documento examinan si los ataques no lineales
tienen éxito. Se organizan por orden de sofisticación algebraica
creciente:

\begin{itemize}
  \item \textbf{Experimentos~07 y~11}: caracterizan la función
    $g_B(y) = y + B\,\sigma(y)$, la función principal que el atacante
    debe invertir.
  \item \textbf{Experimento~08}: establece el umbral exacto bajo el
    cual la clave es recuperable trivialmente ($d = 1$) y confirma
    algebraicamente que los ataques lineales fallan para $d \geq 2$.
  \item \textbf{Experimentos~09 y~12}: miden el coste empírico de la
    recuperación de clave bajo el modelo de texto en claro conocido
    mediante búsqueda exhaustiva sobre $A$.
  \item \textbf{Experimento~10}: presenta un ataque CPA directo que
    requiere únicamente $O(d)$ consultas.
  \item \textbf{Experimento~13}: propone un ataque diferencial por
    columnas que recupera las columnas de $A$ de forma independiente,
    reduciendo el espacio de búsqueda de $|\GL(d,p)|$ a
    $d\,(p^d - 1)$ candidatos.
\end{itemize}

Todos los experimentos emplean la implementación de producción en
\texttt{alaniz/} sin modificación. Las semillas aleatorias están
fijadas para reproducibilidad total.

% ================================================================
\section{Notación del Protocolo}
% ================================================================

La Tabla~\ref{tab:notacion} resume las magnitudes utilizadas a lo
largo de este documento.

\begin{table}[h]
\centering
\caption{Notación empleada en este análisis.}
\label{tab:notacion}
\begin{tabularx}{\linewidth}{>{\ttfamily}lX}
\toprule
\normalfont Símbolo & Significado \\
\midrule
$p$                & Característica del cuerpo (número primo) \\
$d$                & Dimensión de la fibra \\
$G$                & Grafo con $n$ nodos \\
$\Hzero(G, \calF)$ & Espacio de secciones globales (espacio de texto en claro) \\
$s_v \in \Fp^d$    & Texto en claro local en el nodo $v$ \\
$(A_v, B_v)$       & Clave secreta en el nodo $v$; ambas en $\GL(d,\Fp)$ \\
$\sigma$           & Función pública componente a componente:
                     \texttt{inverse} ($x_i \mapsto x_i^{p-2}$)
                     o \texttt{cube} ($x_i \mapsto x_i^3$) \\
$y_v = A_v s_v$    & Vector intermedio secreto en el nodo $v$ \\
$g_B(y)$           & $y + B\,\sigma(y)$: función de cifrado local \\
\bottomrule
\end{tabularx}
\end{table}

El descifrado resuelve $y + B_r\,\sigma(y) = c_r$ en el nodo raíz,
propaga mediante morfismos de restricción y usa filtrado cohomológico
para seleccionar la sección global válida.

% ================================================================
\section{Experimento 07 --- Tabla de Colisiones de la Función de Cifrado Local}
% ================================================================

\subsection{Objetivo}

Cuantificar la sobreyectividad y la multiplicidad de preimagen de
$g_B$ sobre el dominio completo $\Fp^d$ para $d \in \{1, 2, 4\}$ y
$p \in \{17, 23, 29, 47, 59, 89, 131, 251\}$.
Esto aporta evidencia concreta sobre si $g_B$ es inyectiva, lo que
determina la ambigüedad en el descifrado.

\subsection{Metodología}

Para $d \in \{1, 2\}$, se enumeran exactamente todas las $p^d$
entradas. Se construye el mapa de preimagen $\{c \mapsto g_B^{-1}(c)\}$
y se calculan las siguientes métricas por matriz aleatoria
$B \in \GL(d, \Fp)$ (promediadas sobre ocho realizaciones
independientes):
\begin{equation}
  \text{fracción de imagen} = \frac{|\Img(g_B)|}{p^d}, \qquad
  \gamma_B = \sum_{c \in \Fp^d} \!\left(\frac{|g_B^{-1}(c)|}{p^d}\right)^{\!2}.
\end{equation}
Para $d = 4$, donde la enumeración exhaustiva no es viable, se emplea
un sondeo de cumpleaños con $3\lceil\sqrt{p^d}\rceil$ muestras
aleatorias, que garantiza encontrar una colisión con alta probabilidad.
Toda colisión hallada se verifica frente a la condición algebraica
$\Delta y = B\bigl(\sigma(y_2) - \sigma(y_1)\bigr)$.

\subsection{Resultados}

En todos los primos y ambas funciones $\sigma$, la fracción de imagen
se estabiliza en torno a $0{,}53$ para $\sigma = \texttt{inverse}$ y
$0{,}66$ para $\sigma = \texttt{cube}$ en dimensión $d = 1$, y en torno
a $0{,}62$--$0{,}64$ para $d = 2$. La función $g_B$ es, por tanto,
consistentemente no sobreyectiva: aproximadamente un tercio de los
textos cifrados en cada nodo son inalcanzables bajo una clave fija.
Se confirman colisiones en todos los tamaños de parámetros, incluyendo
$d = 4$, $p = 251$.

\subsection{Implicaciones para la Seguridad}

La no sobreyectividad tiene dos consecuencias. Primero, el descifrado
en el nodo raíz admite múltiples preimágenes (multiplicidad máxima
observada de 2--5 para $d = 2$), introduciendo una ambigüedad que
resuelve el filtro cohomológico. Segundo, la probabilidad de colisión
$\gamma_B$ proporciona una estimación empírica directa del parámetro
$\gamma$ del Teorema~7.1 del artículo original.

% ================================================================
\section{Experimento 08 --- Recuperación de Clave en $d=1$ y Diagnóstico de Ataques Lineales}
% ================================================================

\subsection{Objetivo}

Establecer el umbral exacto bajo el cual la clave es recuperable
algebraicamente, y confirmar que la defensa mediante linealización
se mantiene para $d \geq 2$.

\subsection{Recuperación Exacta para $d = 1$}

Para $d = 1$, el cifrado se reduce a una ecuación escalar. Dados dos
pares texto en claro--texto cifrado $(s_1, c_1)$ y $(s_2, c_2)$:

\paragraph{Caso $\sigma = \texttt{inverse}$.}
Multiplicando $c_i = a s_i + b(as_i)^{-1}$ por $as_i$ y restando
las dos ecuaciones se obtiene:
\begin{equation}
  a = \frac{c_1 s_1 - c_2 s_2}{s_1^2 - s_2^2}, \qquad
  b = c_1(as_1) - (as_1)^2.
\end{equation}

\paragraph{Caso $\sigma = \texttt{cube}$.}
Las dos ecuaciones forman un sistema $2\times 2$ de tipo Vandermonde
en $(a,\, ba^3)$, resuelto por la regla de Cramer con determinante
$s_1 s_2^3 - s_2 s_1^3$.

Ambas fórmulas son exactas sobre $\Fp$ y requieren únicamente dos
capturas.

\subsection{Fallo del Sustituto Lineal para $d \geq 2$}

Para $d \in \{2, 4\}$, el experimento construye el sistema lineal
sustituto $c_v = M s_v$ (tratando la clave como una única matriz
$d\times d$) y rastrea su rango y consistencia conforme se acumulan
capturas. Como era de esperar, la no linealidad de $\sigma$ hace el
sistema inconsistente a partir de un número suficiente de pares,
confirmando que la linealización pura falla.

\subsection{Resultados e Implicaciones}

La dimensión $d = 1$ no proporciona \textbf{ninguna seguridad
computacional}: la clave se recupera con 2 pares
texto en claro--texto cifrado usando únicamente aritmética modular.
La dimensión mínima segura es $d \geq 2$.
Sin embargo, el Experimento~10 demuestra que un ataque CPA no lineal
tiene éxito en todas las dimensiones.

% ================================================================
\section{Experimento 09 --- Recuperación Exhaustiva de Clave con Texto en Claro Conocido}
% ================================================================

\subsection{Objetivo}

Medir el coste concreto de recuperar la clave por nodo $(A_v, B_v)$
bajo el modelo de texto en claro conocido mediante búsqueda exhaustiva
sobre $A \in \GL(d, \Fp)$ con resolución lineal de $B$.

\subsection{Metodología}

Para cada candidato $\hat{A}$, se calculan $y_i = \hat{A}\, s_i$ y
$z_i = \sigma(y_i)$ para las primeras $d$ capturas y se resuelve el
sistema lineal
\begin{equation}
  c_i - y_i = B\, z_i, \qquad i = 1, \ldots, d,
\end{equation}
para las $d^2$ entradas de $B$. El candidato $(\hat{A}, \hat{B})$
se valida frente a las $k$ capturas restantes. La campaña cubre
$p \in \{17, 23, 29, 47, 59, 89, 131, 251\}$, $d = 2$, con
$k \in \{4, 8, 12\}$ capturas. Para $|\GL(2, p)| \leq 10^5$ se usa
enumeración exhaustiva; para primos mayores, 5\,000 muestras
aleatorias.

\subsection{Resultados e Implicaciones}

Para $p \leq 29$, con $|\GL(2, p)| \leq 7\times 10^5$, la clave
verdadera se recupera exactamente en todos los ensayos. Para
$p \geq 47$ la probabilidad de muestreo es despreciable
(${\approx}5000/|\GL(2,p)|$). Los falsos positivos desaparecen con
$k$: en $k = 12$ no sobrevive ninguno en ningún ensayo exhaustivo.
La búsqueda exhaustiva sobre $A$ es intratable para los parámetros
estándar propuestos ($d = 4$, $p \approx 2^{31}$), pero el
Experimento~10 hace que este límite sea irrelevante al recuperar la
clave con $O(d)$ consultas de texto en claro elegido con independencia
de $p$.

% ================================================================
\section{Experimento 10 --- Ataque CPA Directo mediante Homogeneidad de Escala}
% ================================================================

\textit{Este es el hallazgo de seguridad principal de este análisis.}

\subsection{Objetivo}

Demostrar un ataque de texto en claro elegido que recupera
$(A_v, B_v)$ en $O(d)$ consultas de cifrado, sin ninguna búsqueda
exhaustiva sobre $\GL(d, \Fp)$.

\subsection{Base Matemática}
\label{sub:escala}

Ambas funciones no lineales soportadas satisfacen una
\textbf{ley de escala componente a componente}:
\begin{equation}\label{eq:homogeneidad}
  \sigma(\lambda\, y)_i = \lambda^e\, \sigma(y)_i
  \qquad \forall\, y \in \Fp^d,\; \lambda \in \Fp^*,
\end{equation}
con $e = 3$ para $\sigma = \texttt{cube}$ y
$e \equiv p-2 \pmod{p-1}$ para $\sigma = \texttt{inverse}$.

Dado que $B$ actúa linealmente sobre $\sigma(y)$, la
ecuación~\eqref{eq:homogeneidad} se extiende al texto cifrado
completo. Para textos en claro $s = t\,e_j$ (múltiplos escalares del
$j$-ésimo vector de la base canónica del espacio de coeficientes de
$H^0$), el texto cifrado local satisface
\begin{equation}\label{eq:par}
  c(t\,e_j) = t\,a_j + t^e\,B\,\sigma(a_j),
\end{equation}
donde $a_j = A\,e_j$ es la $j$-ésima columna de $A$.
Consultando $s_1 = e_j$ ($t = 1$) y $s_2 = \lambda\,e_j$
($t = \lambda$) y escribiendo $w_j = B\,\sigma(a_j)$:
\begin{equation}\label{eq:sistema2x2}
  \begin{pmatrix} 1 & 1 \\ \lambda & \lambda^e \end{pmatrix}
  \begin{pmatrix} a_j \\ w_j \end{pmatrix}
  = \begin{pmatrix} c_1 \\ c_2 \end{pmatrix}.
\end{equation}
El determinante es $\Delta = \lambda^e - \lambda$. Eligiendo cualquier
$\lambda$ con $\Delta \neq 0$ (lo que es posible para todo $p > 3$)
se obtiene la solución en \emph{forma cerrada}:
\begin{equation}\label{eq:solucion}
  a_j = \frac{\lambda^e c_1 - c_2}{\Delta}, \qquad
  w_j = \frac{-\lambda c_1 + c_2}{\Delta}.
\end{equation}

Repitiendo para $j = 0, \ldots, d-1$ se recupera la matriz de clave
efectiva $A_{\text{ef}} := A\,P$ en $2d$ consultas, donde $P$ es el
mapa coeficiente--sección, computable públicamente a partir del haz.
La clave $A$ en coordenadas de sección se recupera como:
\begin{equation}
  A = A_{\text{ef}} \cdot P^{-1}.
\end{equation}

Con $A_{\text{ef}}$ conocida, $B$ se recupera recopilando $d$
consultas adicionales con $z^{(i)} = \sigma(A_{\text{ef}}\,s^{(i)})$
linealmente independientes y resolviendo:
\begin{equation}\label{eq:recuperacion-B}
  B = W\,Z^{-1}, \qquad
  W_i = c^{(i)} - A_{\text{ef}}\,s^{(i)}, \quad
  Z = \bigl[z^{(1)} \cdots z^{(d)}\bigr].
\end{equation}
Número total de consultas: $2d + d = 3d$ en el peor caso, o
$2d + 1$ cuando las consultas de base ya proporcionan $Z$ invertible.

\subsection{Validación Experimental}

El ataque se ejecuta frente a todos los primos
$p \in \{17, 23, 29, 47, 59, 89, 131, 251\}$, ambas funciones
$\sigma$ y dimensiones $d \in \{2, 4\}$. En todos los casos:

\begin{itemize}
  \item $A_{\text{ef}}$ se recupera exactamente
    (\texttt{A\_eff\_match = True}).
  \item $A$ en coordenadas de sección se recupera exactamente
    (\texttt{A\_stalk\_match = True}).
  \item $B$ se recupera exactamente (\texttt{B\_match = True}).
  \item La clave recuperada supera 20 consultas de validación
    independientes (\texttt{valid = True}).
\end{itemize}

Para $d = 4$, el número total de consultas de cifrado empleadas es
$12$, con independencia de $p$.

\subsection{Implicaciones para la Seguridad}
\label{sub:implicaciones}

Este resultado tiene las siguientes consecuencias para las
afirmaciones de seguridad del artículo original:

\begin{enumerate}

  \item \textbf{La afirmación de resistencia CPA queda refutada para
    todos los conjuntos de parámetros.} El ataque no requiere
    linealización y no es bloqueado por la no linealidad de $\sigma$.
    Al contrario, la homogeneidad de escala de $\sigma$ \emph{es}
    la propiedad algebraica que hace posible el ataque. La afirmación
    ``la linealización falla porque $B_v\,\sigma(A_v s_v)$ mezcla
    clave y mensaje de forma no lineal'' sigue siendo verdadera, pero
    no cubre este vector de ataque.

  \item \textbf{La reducción de dureza NL-SMIP no se ve afectada.}
    La reducción desde MQ establece que recuperar la clave a partir
    de un sistema polinómico con coeficientes secretos es NP-duro.
    El ataque de escala no resuelve NL-SMIP; lo evita completamente
    explotando una propiedad estructural de $\sigma$ que no captura
    la formulación polinómica.

  \item \textbf{Las afirmaciones post-cuánticas no se ven
    sustancialmente afectadas.} El ataque de escala es puramente
    algebraico y se aplica con la misma complejidad $O(d)$ de forma
    clásica y cuántica. No explota ninguna aceleración cuántica;
    simplemente es más rápido que cualquier búsqueda exhaustiva.

  \item \textbf{El ataque se extiende simultáneamente a todos los
    nodos.} La clave de cada nodo $(A_v, B_v)$ es independiente y
    el ataque es local al nodo. Consultar los $n$ nodos requiere
    $3dn$ consultas de cifrado. Para un camino de 6 nodos con
    $d = 4$, esto es 72 consultas.

  \item \textbf{El ataque es reparable.} Cualquier $\sigma$ que
    satisfaga $\sigma(\lambda y) \neq \lambda^e \sigma(y)$ para todo
    $e$ derrotaría este ataque concreto. Extender el análisis de
    selección de $\sigma$ del Experimento~04 para incluir la
    homogeneidad de escala es la dirección natural para una
    corrección.

\end{enumerate}

% ================================================================
\section{Experimento 11 --- Estudio Estadístico de Colisiones de $g_B$}
% ================================================================

\subsection{Objetivo}

Obtener estimaciones estadísticamente robustas de las propiedades de
colisión de $g_B$ en todo el espacio de parámetros, ampliando la
verificación puntual del Experimento~07 a grandes muestras de
matrices $B$.

\subsection{Metodología}

Para $d = 2$ y todos los primos, la enumeración exacta sobre $p^d$
entradas se repite para 12--24 realizaciones independientes de
$B \in \GL(d,\Fp)$. Para $d = 4$, se muestrean 30\,000 entradas
aleatorias por $B$. Por configuración se registran:
\begin{itemize}
  \item Fracción de imagen media $\pm$ desviación típica sobre la
    muestra de $B$.
  \item Multiplicidad media de preimagen: $p^d / |\Img(g_B)|$.
  \item Ejemplos de las peores matrices $B$ (fracción de imagen
    mínima).
  \item Histograma de multiplicidad agregado.
\end{itemize}

\subsection{Resultados e Implicaciones}

La fracción de imagen para $\sigma = \texttt{inverse}$ se estabiliza
en torno a $0{,}63$ para $d = 2$ con varianza muy baja
($\sigma_{\text{std}} < 0{,}01$ para $p \geq 59$), confirmando que
la no sobreyectividad es una propiedad estable de $\sigma$ y no un
artefacto de matrices $B$ desfavorables. La multiplicidad máxima de
preimagen permanece acotada por debajo de 7 para $d = 2$,
$p \leq 251$, en consonancia con el Teorema~7.1. Junto con el
Experimento~07, este experimento proporciona estimaciones empíricas
de $\gamma \approx 0{,}37$ y $L \approx 1{,}59$ para la configuración
$\sigma = \texttt{inverse}$, $d = 2$.

% ================================================================
\section{Experimento 12 --- Diagnóstico de Recuperación de Clave y Análisis de Capturas Mínimas}
% ================================================================

\subsection{Objetivo}

Medir el número mínimo de pares de texto en claro conocido necesarios
para la recuperación única de clave en $d = 2$ (exhaustivo) y
$d = 3$ (muestreado), y registrar con qué rapidez se eliminan los
falsos positivos conforme se acumulan capturas.

\subsection{Metodología}

Para cada candidato $\hat{A}$ (exhaustivo para $d = 2$, $p = 17$;
60\,000 muestras aleatorias para $d = 3$), $B$ se resuelve a partir
de las primeras $d$ capturas y el candidato resultante se contrasta
frente a un prefijo creciente de hasta 20 capturas. Se registra el
\emph{índice de primer fallo}---el menor $k$ en el que el candidato
falla---. El mínimo $k$ para el que exactamente un candidato (la
clave verdadera) sobrevive se denota $k_{\min}$.

\subsection{Resultados e Implicaciones}

Para $d = 2$, $p = 17$: el número de candidatos supervivientes cae de
miles en $k = d = 2$ a pocos en $k = 6$, y exactamente a 1 (sin
falsos positivos) en $k \approx 10$--$15$, dependiendo de $\sigma$.
Para $d = 3$ en modo muestreado, la recuperación única se logra en
$k \approx 12$--$18$ capturas. Cada captura aporta aproximadamente
$d^2 \log_2 p$ bits de información sobre la clave, en consonancia
con la estructura del sistema polinómico. Estos límites son
principalmente relevantes para configuraciones de prueba; en los
parámetros estándar propuestos, la búsqueda exhaustiva sobre $A$ es
intratable (Experimento~09).

% ================================================================
\section{Experimento 13 --- Ataque Diferencial por Columnas}
% ================================================================

\subsection{Objetivo}

Comprobar si las columnas de $A$ pueden recuperarse \emph{de forma
independiente} mediante consultas diferenciales de texto en claro
elegido a lo largo de los ejes coordenados, y medir la reducción
resultante en la complejidad del ataque.

\subsection{Base Matemática}

Para textos en claro $s = t\,e_j$, las diferencias consecutivas de
texto cifrado son:
\begin{equation}\label{eq:diferencia}
  \Delta c_t = c\bigl((t+1)e_j\bigr) - c\bigl(t\,e_j\bigr)
  = a_j + B\bigl(\sigma\bigl((t+1)a_j\bigr) - \sigma(t\,a_j)\bigr).
\end{equation}
Para una columna candidata $\hat{a}_j$, esto es una ecuación lineal
en $B$. Recopilando $d$ diferencias y denotando
$\Delta\sigma_t = \sigma((t+1)\hat{a}_j) - \sigma(t\hat{a}_j)$:
\begin{equation}
  \Delta c_t - \hat{a}_j = B\,\Delta\sigma_t.
\end{equation}
Si la matriz $[\Delta\sigma_0 \cdots \Delta\sigma_{d-1}]$ es
invertible, $B$ queda determinada unívocamente.
Cada columna de $A$ es recuperable buscando entre los $p^d - 1$
candidatos no nulos de $\hat{a}_j$, reduciendo el espacio de
búsqueda total de $|\GL(d,p)| \approx p^{d^2}$ a $d\,(p^d - 1)$.

\subsection{Resultados}

\begin{description}
  \item[$d = 2$, $p = 17$:] $p^2 - 1 = 288$ candidatos por eje
    \emph{frente a} $|\GL(2,17)| = 69\,360$ matrices completas.
    Recuperación única en 4--6 diferencias en todas las
    configuraciones probadas.
  \item[$d = 3$, $p = 17$:] $p^3 - 1 = 4\,912$ candidatos
    \emph{frente a} $|\GL(3,17)| \approx 1{,}4\times 10^7$.
    Recuperación única en 4--7 diferencias.
\end{description}

\paragraph{Limitación.} Si la columna $a_j$ tiene un cero en la
componente $i$, entonces $\sigma(t\,a_j)_i = 0$ para todo $t$,
haciendo que la $i$-ésima fila de $\Delta\sigma$ sea idénticamente
nula. Las entradas correspondientes de $B$ permanecen indeterminadas
y el ataque produce falsos negativos en esas configuraciones.

\subsection{Implicaciones para la Seguridad}

El ataque diferencial por columnas reduce la búsqueda por columna de
$p^{d^2}$ a $p^d - 1$ candidatos---una mejora exponencial para
$d \geq 2$. Queda dominado por el ataque de $O(d)$ consultas del
Experimento~10 en términos prácticos, pero proporciona una vía de
recuperación independiente útil cuando el acceso directo de consultas
está restringido.

% ================================================================
\section{Discusión}
% ================================================================

\subsection{La Homogeneidad de Escala como Vulnerabilidad Estructural}

La causa raíz del ataque CPA del Experimento~10 es la interacción
entre tres decisiones de diseño:
\begin{enumerate}
  \item $\sigma$ se aplica \emph{componente a componente}, lo que
    permite $\sigma(\lambda y)_i = \lambda^e \sigma(y)_i$.
  \item $B$ actúa \emph{linealmente} sobre $\sigma(y)$, de modo que
    $B\,\sigma(\lambda y) = \lambda^e\,B\,\sigma(y)$.
  \item El espacio de texto en claro contiene \emph{múltiplos
    escalares} de cualquier vector, permitiendo al atacante consultar
    $(s, \lambda s)$ y desacoplar $y = As$ de
    $w = B\sigma(As)$ mediante un sistema lineal $2\times 2$.
\end{enumerate}
Cualquiera de las dos primeras propiedades por sí sola no sería
suficiente para el ataque. Una corrección debe, por tanto, eliminar
al menos una de ellas. La opción más natural es sustituir $\sigma$
por una función que carezca de escala componente a componente,
preservando al mismo tiempo la biyectividad y la compatibilidad con
el descifrado.

\subsection{Relación con la Prueba de Seguridad Original}

El argumento de resistencia CPA del artículo original (Sección~5)
considera diseños en los que $\sigma$ se aplica antes o con
independencia de la transformación secreta, demostrando que en esos
casos la clave puede recuperarse linealmente. El diseño final sitúa
$\sigma$ \emph{después} de la transformación lineal secreta $A_v$,
lo que sí derrota la linealización. El Experimento~10 muestra que
derrotar la linealización es necesario pero no suficiente para la
seguridad CPA: la homogeneidad de escala proporciona un ataque
alternativo que no requiere linealizar el sistema.

La prueba de dureza NL-SMIP (Lema~6.2 del artículo original)
establece que recuperar la clave a partir de un sistema polinómico
con coeficientes secretos es NP-duro. El ataque de escala modela
un escenario diferente---en el que el atacante controla los textos
en claro y explota la estructura algebraica de $\sigma$---. Ambos
modelos no son equivalentes, y la NP-dureza de NL-SMIP no implica
seguridad CPA.

\subsection{Propiedades de Colisión y Corrección del Descifrado}

Los Experimentos~07 y~11 establecen conjuntamente que $g_B$ es no
sobreyectiva con fracción de imagen $\approx 0{,}63$ para $d = 2$.
Esto proporciona estimaciones empíricas de $\gamma \approx 0{,}37$
y $L \approx 1{,}59$ para la cota de corrección del Teorema~7.1
($\Pr[\text{ambigüedad}] \leq (L-1)\,\gamma^{n-1}$). Para grafos con
$n \geq 6$ nodos, la ambigüedad en el descifrado es despreciable en
la práctica.

% ================================================================
\section{Conclusión}
% ================================================================

Este análisis presenta siete experimentos de seguridad reproducibles
que sondean el Cifrado Alaniz desde múltiples perspectivas
adversariales. Las conclusiones principales son:

\begin{enumerate}
  \item \textbf{La dimensión $d = 1$ es trivialmente insegura}: la
    clave se recupera algebraicamente a partir de 2 pares
    texto en claro--texto cifrado (Experimento~08).

  \item \textbf{La función local $g_B$ es consistentemente no
    sobreyectiva} (fracción de imagen $\approx 0{,}63$) con
    multiplicidad de preimagen acotada, validando el análisis de
    unicidad del descifrado (Experimentos~07 y~11).

  \item \textbf{La recuperación exhaustiva de clave con texto en
    claro conocido} es viable para parámetros de juguete ($d = 2$,
    $p \leq 29$) y requiere aproximadamente 10--15 capturas para
    recuperación única; es intratable para los parámetros estándar
    propuestos (Experimentos~09 y~12).

  \item \textbf{Existe un ataque CPA directo para todos los
    parámetros}: la homogeneidad de escala
    $\sigma(\lambda y) = \lambda^e\,\sigma(y)$ permite recuperar la
    clave por nodo $(A_v, B_v)$ con $O(d)$ consultas de texto en
    claro elegido, con independencia de $p$ (Experimento~10). Este
    es el hallazgo de seguridad principal de este análisis.

  \item \textbf{Un ataque diferencial por columnas} reduce el espacio
    de búsqueda por columna de $|\GL(d,p)|$ a $d\,(p^d - 1)$
    candidatos (Experimento~13).
\end{enumerate}

El ataque de escala del Experimento~10 es una vulnerabilidad
significativa que requiere atención antes de que el esquema pueda
considerarse seguro conforme a las definiciones CPA estándar. La
dirección natural para una corrección es sustituir la $\sigma$
homogénea de escala componente a componente por una función que mezcle
componentes a través de la fibra $d$-dimensional, preservando las
propiedades algebraicas necesarias para el descifrado eficiente.

% ================================================================
\section*{Tabla Resumen de Experimentos}
\addcontentsline{toc}{section}{Tabla Resumen de Experimentos}
% ================================================================

\begin{table}[h]
\centering
\caption{Resumen de los Experimentos 07--13.}
\label{tab:resumen}
\renewcommand{\arraystretch}{1.35}
\begin{tabularx}{\linewidth}{c l l l X}
\toprule
Exp & Modelo & Ataque / Análisis & Complejidad & Resultado Principal \\
\midrule
07 & ---  & Caracterización de colisiones de $g_B$
           & $O(p^d \cdot b)$
           & Fracción de imagen $\approx 0{,}63$; colisiones universales \\
08 & TCC  & Recuperación algebraica de clave
           & $O(1)$ para $d{=}1$
           & $d{=}1$ trivialmente inseguro; linealización falla en $d{\geq}2$ \\
09 & TCC  & Búsqueda exhaustiva sobre $A$ + $B$ lineal
           & $O(|\GL(d,p)|)$
           & Viable para $p \leq 29$; intratable en parámetros estándar \\
\textbf{10} & \textbf{TCE}
           & \textbf{Ataque de homogeneidad de escala}
           & $\bm{O(d)}$ \textbf{consultas}
           & \textbf{Recuperación completa en todos los parámetros; CPA refutado} \\
11 & ---  & Estudio estadístico de multiplicidad de $g_B$
           & $O(p^d \cdot b)$
           & Fracción de imagen estable; familias de peor $B$ identificadas \\
12 & TCC  & Diagnóstico de capturas mínimas
           & $O(|\GL(d,p)| \cdot k)$
           & Recuperación única en $\approx 10$--$15$ capturas para $d{=}2$ \\
13 & TCE  & Ataque diferencial por columnas
           & $O(d(p^d{-}1))$
           & Espacio de búsqueda reducido de $p^{d^2}$ a $p^d{-}1$ por columna \\
\bottomrule
\multicolumn{5}{l}{\small TCC = Texto en Claro Conocido; \quad TCE = Texto en Claro Elegido.}
\end{tabularx}
\end{table}

\bigskip
\noindent
\textit{Los Experimentos 07--13 están implementados en
\texttt{experiments/07\_collision\_table.py} hasta
\texttt{experiments/13\_differential\_column\_attack.py}.
Todas las semillas aleatorias están fijadas; los resultados son
completamente reproducibles.}

\end{document}
